%====================================
% KAPITEL 1: EINLEITUNG
%====================================
\section{Einleitung}

Bereits im Jahr 2003 veröffentlichten DeLone und McLean mit ihrem aktualisierten Information-Systems-Success-Modell einen der bis heute einflussreichsten Ansätze zur Bewertung des Erfolgs von Informationssystemen.\footcite[Vgl.][S.~9--30]{delone2003} Das Modell beschreibt den Erfolg technologischer Systeme anhand mehrerer miteinander verknüpfter Dimensionen und gilt als zentraler Referenzrahmen der Informationssystemforschung. Trotz dieses etablierten Verständnisses davon, was den Erfolg eines Systems ausmacht, bleiben ERP-Implementierungen in der Praxis jedoch häufig hinter den Erwartungen zurück. Aktuelle systematische Literaturanalysen zeigen, dass zahlreiche Unternehmen die angestrebten Potenziale ihrer ERP-Systeme selbst nach der Einführung nicht vollständig realisieren.\footcite[Vgl.][S.~2--3]{butarbutar2023}

Diese Befunde verdeutlichen, dass technologische Exzellenz allein nicht ausreicht. Bereits im Jahr 2001 systematisierten Nah, Lau und Kuang die elf kritischen Erfolgsfaktoren für ERP-Projekte und identifizierten dabei das Change Management explizit als eigenständigen Faktor.\footcite[Vgl.][S.~285]{nah2001} Die Forschung zeigt, dass Change Management und die Unterstützung durch das Top-Management zu den zentralen Erfolgsfaktoren von ERP-Implementierungen zählen.\footcites[Vgl.][S.~285]{nah2001}[][S.~257--278]{somers2004} Diese Erfolgsproblematik ist nicht neu: Schon Holland und Light wiesen darauf hin, dass rein technologische Ansätze ohne Berücksichtigung der organisatorischen Passung oft defizitär bleiben.\footcite[Vgl.][S.~30--36]{holland1999} Systematische Literaturübersichten bestätigen zudem, dass viele Unternehmen auch in der Post-Implementierungsphase Schwierigkeiten haben, die angestrebten Potenziale voll auszuschöpfen.\footcite[Vgl.][Art.-Nr.~2264001]{butarbutar2023}

Um diesen Herausforderungen zu begegnen, bietet das Change Management die notwendige theoretische sowie praktische Perspektive. Während Armenakis und Bedeian die Grundlagen für den Umgang mit organisationalem Wandel schufen,\footcite[Vgl.][S.~293--315]{armenakis1999} verdeutlicht die Forschung von Hubbart, dass insbesondere die individuelle Abwehrhaltung gegenüber Veränderungen eine erhebliche Hürde bleibt.\footcite[Vgl.][S.~1~f.]{hubbart2023} Den Übergang zwischen allgemeinen Strategien und dem spezifischen ERP-Kontext bildet der Ansatz von Aladwani, der betont, dass der Projekterfolg wesentlich von einem strukturierten, prozessorientierten Umgang mit Nutzerwiderstand abhängt.\footcite[Vgl.][S.~266--275]{aladwani2001} Als zentrale Wirkmechanismen gelten dabei die Hebel Kommunikation, Führung und Schulung.\footcites[Vgl.][S.~285]{nah2001}[][Art.-Nr.~2264001]{butarbutar2023} Obwohl die Relevanz dieser Faktoren in der Forschung breit belegt ist, fokussiert der Gro\ss{}teil der Studien bislang Gro\ss{}unternehmen und Konzerne.\footcite[Vgl.][Art.-Nr.~2264001]{butarbutar2023}

Hinzu kommt jedoch die volkswirtschaftliche Relevanz: Der Mittelstand bildet das Rückgrat der deutschen Wirtschaft, muss die digitale Transformation jedoch häufig mit deutlich begrenzteren Ressourcen bewältigen als Gro\ss{}konzerne. Aufgrund dieser spezifischen Restriktionen und flacheren Hierarchien bleiben mittelständische Unternehmen in der Forschung oft unterrepräsentiert, weshalb es an validierten Handlungsempfehlungen für diesen speziellen Kontext fehlt. Zentral ist die Beantwortung der Frage: Welche Change-Management-Faktoren beeinflussen den Erfolg von ERP-Implementierungen in mittelständischen Unternehmen?

Zur Beantwortung dieser Frage wird der Erfolg auf Basis des DeLone-\&-McLean-IS-Success-Modells operationalisiert.\footcite[Vgl.][S.~9--30]{delone2003} Die einschlägige Forschungsliteratur ab dem Jahr 2015 wird nach dem methodischen Vorgehen von Webster und Watson systematisch ausgewertet.\footcite[Vgl.][S.~xiii--xxiii]{webster2002} Ein besonderer Fokus liegt dabei auf der Analyse der Hebel Kommunikation, Führung und Schulung, die im Rahmen des ADKAR-Modells als zentrale Wirkdimensionen verortet werden,\footcite[Vgl.][S.~124--132]{galli2018} um praxisbezogene Empfehlungen für die Gestaltung von ERP-Projekten abzuleiten. Die Arbeit gliedert sich zunächst in einen theoretischen Grundlagenteil (Kapitel 2), gefolgt von der Darstellung des methodischen Vorgehens (Kapitel 3). In Kapitel 4 werden die Ergebnisse der systematischen Literaturanalyse präsentiert, bevor in Kapitel 5 die daraus resultierenden Handlungsempfehlungen für den Mittelstand abgeleitet werden. Die Arbeit schlie\ss{}t mit einem Fazit und einem Ausblick (Kapitel 6).
